\documentclass[conference]{IEEEtran}
\IEEEoverridecommandlockouts
% The preceding line is only needed to identify funding in the first footnote. If that is unneeded, please comment it out.
\usepackage{cite}
\usepackage{amsmath,amssymb,amsfonts}
\usepackage{algorithmic}
\usepackage{graphicx}
\usepackage{textcomp}
\usepackage{xcolor}
\usepackage{booktabs}
\usepackage{siunitx}
\usepackage[T1]{fontenc}
\usepackage{eurosym}
\usepackage[caption=false]{subfig}
\usepackage{listings}
\usepackage[T1]{fontenc}
\usepackage[utf8]{inputenc}
% \usepackage[portuges]{babel} // para a língua Portuguesa
% \selectlanguage{portuges} // para colocar em Português o título do abstract e das referências
\usepackage{orcidlink} % to typeset URLs, URIs, and DOIs
\def\UrlFont{\rmfamily} % to typeset the Euro symbol
\DeclareSIUnit{\sieuro}{\mbox{\euro}}
\def\BibTeX{{\rm B\kern-.05em{\sc i\kern-.025em b}\kern-.08em
    T\kern-.1667em\lower.7ex\hbox{E}\kern-.125emX}}
\definecolor{darkblue}{rgb}{0.0,0.0,0.6}
\lstset{
  basicstyle=\ttfamily,
  basicstyle=\footnotesize,
  columns=fullflexible,
  showstringspaces=false,
  commentstyle=\color{gray}\upshape
}

\lstdefinelanguage{XML}
{
  morestring=[b]",
  morestring=[s]{>}{<},
  morecomment=[s]{<?}{?>},
  stringstyle=\color{black},
  identifierstyle=\color{darkblue},
  keywordstyle=\color{cyan},
  morekeywords={name}% list your attributes here
}

\begin{document}

\title{Conference Paper Title %*\\
%{\footnotesize \textsuperscript{*}Note: Sub-titles are not captured in Xplore and should not be used}
%\thanks{Identify applicable funding agency here. If none, delete this.}
}

\author{\IEEEauthorblockN{Given Name Surname}
\IEEEauthorblockA{\textit{dept. name of organization (of Aff.)} \\
\textit{name of organization (of Aff.)}\\
City, Country \\
email address or ORCID}
\and
\IEEEauthorblockN{Given Name Surname}
\IEEEauthorblockA{\textit{dept. name of organization (of Aff.)} \\
\textit{name of organization (of Aff.)}\\
City, Country \\
email address or ORCID}
\and
\IEEEauthorblockN{Given Name Surname}
\IEEEauthorblockA{\textit{dept. name of organization (of Aff.)} \\
\textit{name of organization (of Aff.)}\\
City, Country \\
email address or ORCID}
}

\maketitle

\begin{abstract}
This document is a model and instructions for \LaTeX.
This and the IEEEtran.cls file define the components of your paper [title, text, heads, etc.]. *CRITICAL: Do Not Use Symbols, Special Characters, Footnotes, 
or Math in Paper Title or Abstract.
\end{abstract}

\begin{IEEEkeywords}
component, formatting, style, styling, insert
\end{IEEEkeywords}

\section{Introduction}
This document is a model and instructions for \LaTeX.
Please observe the conference page limits. 
Please substitute the instructions by your own text.

\subsection{Background}
The IEEEtran class file is used to format your paper and style the text. All margins, column widths, line spaces, and text fonts are prescribed; please do not alter them. You may note peculiarities. For example, the head margin measures proportionately more than is customary. This measurement and others are deliberate, using specifications that anticipate your paper as one part of the entire proceedings, and not as an independent document. 

\subsection{Objectives}
Define abbreviations and acronyms the first time they are used in the text, even after they have been defined in the abstract. For example, Robotic Operation System (ROS), Personal Computer (PC), Pulse Width Modulation (PWM), etc.

It is mandatory to adopt the International System of Units (SI) \cite{siunitx_2023}. Use the \texttt{\small siunitx} package for this purpose. For example, \SI{2.3}{\metre\per\second}, \SI{5}{\mega\hertz}, \SI{56}{\percent}, etc. Digit separation is automatic with siunitx: \num{123456789.098765432}.

\subsection{Structure of the Report}
Number equations consecutively. Italicize Roman symbols for quantities and variables, but not Greek symbols. Use a long dash rather than a hyphen for a minus sign. Always refer to and explain equations in sentences. For example, Equation \eqref{eq:force} defines force $f$ as the product of the mass $m$ by the acceleration $a$.
\begin{equation}
a+b=\gamma\label{eq:sum}
\end{equation}
\begin{equation}
f=ma\label{eq:force}
\end{equation} 
Be sure that the symbols in your equation have been defined before or immediately following the equation. Use 'Equation \eqref{eq:force}' to refer to each equation in the text.

\section{Problem and Proposed Solution}

\subsection{Problem}
Place figures and tables at the top and bottom of columns. Avoid placing them in the middle of columns. Large figures and tables may span across both columns. Figure captions should be below the figures; table heads should appear above the tables. Insert figures and tables after they are cited in the text. Use ``Figure ~\ref{fig:example}'' to refer to the figure and `Table \ref{tab:example}'' to the table in a sentence.

\begin{table}[htbp]
\caption{Table Type Styles}
\begin{center}
\begin{tabular}{lccc}
\toprule
\textbf{Table}&\multicolumn{3}{c}{\textbf{Table Column Head}} \\
\cmidrule{2-4} 
\textbf{Head} & \textbf{\textit{Table column subhead}}& \textbf{\textit{Subhead}}& \textbf{\textit{Subhead}} \\
\midrule
copy   & More about copy$^{\mathrm{a}}$ & &  \\
remove & To remove ... & &  \\
edit   & To edit ... & &  \\
\bottomrule
\multicolumn{4}{l}{$^{\mathrm{a}}$Sample of a Table footnote.}
\end{tabular}
\label{tab:example}
\end{center}
\end{table}

\begin{figure}[htbp]
\centerline{\includegraphics{fig1.png}}
\caption{Example of a figure caption.}
\label{fig:example}
\end{figure}
\subsection{Solution}
Here are some citation examples: (i) an article \cite{Duarte2020}, a chapter \cite{Andersen2012}, a paper \cite{Thorleuchter2010} and a website \cite{siunitx_2023}.

\subsubsection{Method}

\subsubsection{Technologies}
ROS + OpenCV

\subsection{Sumary}

\section{Development}

Code should be placed inside the listtings environment and referred. For example, Listing \ref{lst:urdf_file} illustrates a partial URDF listing. 
\begin{lstlisting}[language=xml, caption=URDF model of the robot, label=lst:urdf_file]
<robot xmlns:xacro="http://ros.org/wiki/xacro" name="sailboat">
    <link name="sailboat/base_link">
    </link>
    <joint name="sailboat/hull_joint" type = "fixed">
        <parent link="sailboat/base_link"/>
        <child link="sailboat/hull_link"/>
        <origin xyz="0 0 0.1" rpy="0 0 0" />
    </joint>
</robot>
\end{lstlisting}

\subsection{Package Description} 

\subsection{Experiments} 

\subsection{Sumary}

\section{Conclusion}
	
\subsection{Discussion of Results} 

\subsection{Future Work} 

% \section*{Acknowledgment}

% The preferred spelling of the word ``acknowledgment'' in America is without an ``e'' after the ``g''. Avoid the stilted expression ``one of us (R. B. G.) thanks $\ldots$''. Instead, try ``R. B. G. thanks$\ldots$''. Put sponsor acknowledgments in the unnumbered footnote on the first page.

\bibliographystyle{IEEEtran}
\enlargethispage{-5.1in}
\bibliography{IEEEabrv,references}

\end{document}
